\documentclass[twocolumn,pra,amsmath,amssymb,superscriptaddress]{revtex4-2}

\usepackage{graphicx}
\usepackage{hyperref}
\usepackage{booktabs}
\usepackage{xcolor}

\begin{document}

\title{The Sphere Hierarchy: From $S^3$ to $S^6$\\
and the $G_2$ Correspondence}

\author{J.~Loutey}
\affiliation{Independent Researcher, Kent, Washington}

\date{April 13, 2026}

\begin{abstract}
The Geometric Vacuum (GeoVac) framework exploits conformal
equivalences between discrete graphs and Laplace--Beltrami operators
on spheres: $S^3$ for the Coulomb problem (Paper~7) and $S^5$ for
the harmonic oscillator via the Bargmann--Segal lattice (Paper~24).
We present the full sphere hierarchy $S^1$ through $S^7$,
establishing three results.
First, the dimensional cascade theorem: the degeneracy of the $L$-th
harmonic on $S^d$ equals the sum of harmonic degeneracies on $S^{d-1}$,
$g_L(S^d) = \sum_{l=0}^{L} g_l(S^{d-1})$, providing a universal
algebraic link between adjacent spheres.
Second, the $G_2/S^6$ correspondence: the dimensions of the
symmetric irreducible representations $(l,0)$ of the exceptional Lie
group $G_2$ are algebraically identical to the $S^6$ spherical
harmonic degeneracies for all $l \geq 0$, and the $G_2$ quadratic
Casimir eigenvalue on these irreps satisfies
$C_2(l,0) = 2\,l(l+5)$, exactly twice the $S^6$ Laplace--Beltrami
eigenvalue.  Only the symmetric $(l,0)$ irreps appear in $L^2(S^6)$,
via the Peter--Weyl theorem on the coset space $S^6 = G_2/\mathrm{SU}(3)$.
Third, a classification argument identifies $S^3$ as the unique
sphere satisfying three independent criteria simultaneously:
Fock-projectability with a constant scaling parameter (Paper~23
rigidity theorem), constructibility from flat packing (Paper~0),
and realization as the momentum space of a physical
three-dimensional system.  These results classify the sphere
hierarchy by representation-theoretic structure and delineate the
boundary between universal angular properties and Coulomb-specific
features of the GeoVac framework.
\end{abstract}

\maketitle

%=================================================================
\section{Introduction}
\label{sec:intro}
%=================================================================

The GeoVac framework constructs discrete graph Hamiltonians whose
spectra converge to the Laplace--Beltrami operators on spheres
$S^d$ embedded in $\mathbb{R}^{d+1}$.  The foundational case is
$S^3$: Fock's 1935 stereographic projection~\cite{Fock1935} maps
the three-dimensional Coulomb problem onto the free Laplacian on
the unit three-sphere, and Paper~7 of this series establishes the
conformal equivalence between the discrete graph Laplacian and the
$S^3$ Laplace--Beltrami operator through 18 symbolic proofs.
Paper~0 derives the angular quantum numbers $(\ell, m)$ and the
per-shell degeneracy $2\ell + 1$ from a packing construction in
two-dimensional flat space.  Paper~24 extends the framework to
$S^5$ via the Bargmann--Segal lattice, encoding the
three-dimensional isotropic harmonic oscillator on the holomorphic
sector of the Hardy space $H^2(S^5)$.

These constructions raise a natural question: what is the full
sphere hierarchy?  For each $S^d$, the Laplace--Beltrami operator
has a well-defined spectrum with eigenvalues $\lambda_l = -l(l+d-1)$
and degeneracies $g_l(S^d)$ determined by the harmonic polynomial
spaces.  Which of these spheres admit representation-theoretic
structures analogous to the $\mathrm{SO}(4)/S^3$ Fock
correspondence or the $\mathrm{SU}(3)/S^5$ Bargmann--Segal
correspondence?

This paper presents the hierarchy $S^1$ through $S^7$, proves a
universal dimensional cascade connecting adjacent spheres,
identifies a previously uncharacterized correspondence between
$S^6$ and the exceptional Lie group $G_2$, and establishes $S^3$
as the unique sphere satisfying three independent criteria required
by the GeoVac framework.

%=================================================================
\section{The Sphere Hierarchy}
\label{sec:hierarchy}
%=================================================================

\subsection{Eigenvalues and degeneracies}

For each dimension $d \geq 1$, the Laplace--Beltrami operator
$\Delta_{S^d}$ on the unit sphere $S^d \subset \mathbb{R}^{d+1}$
has eigenvalues
\begin{equation}
\lambda_l(S^d) = -l(l + d - 1), \quad l = 0, 1, 2, \ldots
\label{eq:eigenvalues}
\end{equation}
with degeneracies given by the dimension of the space of degree-$l$
harmonic polynomials in $d+1$ variables:
\begin{equation}
g_l(S^d) = \binom{l+d}{d} - \binom{l+d-2}{d}, \quad l \geq 2,
\label{eq:degeneracy}
\end{equation}
with $g_0 = 1$ and $g_1 = d+1$.  Equivalently,
\begin{equation}
g_l(S^d) = \frac{2l + d - 1}{l + d - 1}\binom{l + d - 2}{d - 1}.
\label{eq:degeneracy_alt}
\end{equation}

Table~\ref{tab:hierarchy} presents the eigenvalues, degeneracies,
symmetry groups, and physical interpretations for $S^1$ through
$S^7$.  The entries for $S^2$, $S^3$, and $S^5$ correspond to
established results in the GeoVac framework; the $S^6$ entry is
the new result of this paper.

\begin{table*}
\caption{The sphere hierarchy $S^1$ through $S^7$.
Eigenvalue: $\lambda_l = -l(l+d-1)$.
Degeneracy: $g_l(S^d)$.
Symmetry: the isometry group of $S^d$ and the coset
representation.
Physical identification: the quantum system (if any) whose
state space is encoded by the sphere.
\label{tab:hierarchy}}
\begin{ruledtabular}
\begin{tabular}{clllll}
$d$ & $\lambda_l$ & $g_l$ & Isometry & Coset & Physical identification \\
\hline
1 & $-l^2$      & $2$ ($l \geq 1$), $1$ ($l=0$) & $\mathrm{O}(2)$ & $S^1$ & Fourier modes; particle on a ring \\
2 & $-l(l+1)$   & $2l+1$                         & $\mathrm{SO}(3)$ & $\mathrm{SO}(3)/\mathrm{SO}(2)$ & Angular momentum; packing cross-section (Paper~0) \\
3 & $-l(l+2)$   & $(l+1)^2$                      & $\mathrm{SO}(4)$ & $\mathrm{SO}(4)/\mathrm{SO}(3)$ & 3D Coulomb / hydrogen (Fock projection, Paper~7) \\
4 & $-l(l+3)$   & $\frac{(l+1)(l+2)(2l+3)}{6}$   & $\mathrm{SO}(5)$ & $\mathrm{SO}(5)/\mathrm{SO}(4)$ & No standard identification \\
5 & $-l(l+4)$   & $\frac{(l+1)(l+2)^2(l+3)}{12}$\footnotemark[1]    & $\mathrm{SO}(6)$ & $\mathrm{SO}(6)/\mathrm{SO}(5)$ & 3D HO, holomorphic sector (Paper~24) \\
6 & $-l(l+5)$   & $\frac{(l+1)(l+2)(l+3)(l+4)(2l+5)}{120}$ & $G_2 \subset \mathrm{SO}(7)$ & $G_2/\mathrm{SU}(3)$ & $G_2$ correspondence (this paper) \\
7 & $-l(l+6)$   & $\frac{(l+1)(l+2)(l+3)^2(l+4)(l+5)}{360}$\footnotemark[2]  & $\mathrm{SO}(8)$ & $\mathrm{SO}(8)/\mathrm{SO}(7)$ & 4D HO, holomorphic sector \\
\end{tabular}
\end{ruledtabular}
\footnotetext[1]{Full $S^5$ harmonic degeneracy
$(l+1)(l+2)^2(l+3)/12$.  The 3D~HO shell
degeneracy $(l+1)(l+2)/2$ is the $\mathrm{SU}(3)$ symmetric irrep
$(l,0)$ dimension, a proper subset selected by the Bargmann--Segal
holomorphic projection (Paper~24).  The full $S^5$ harmonic space
is larger by a factor that grows with $l$.}
\footnotetext[2]{Full $S^7$ harmonic degeneracy.
The 4D~HO shell degeneracy $\binom{l+3}{3}$ is the
$\mathrm{SU}(4)$ symmetric irrep $(l,0,0)$ dimension,
a holomorphic subsector analogous to the $S^5$/3D~HO case.}
\end{table*}


\subsection{The dimensional cascade theorem}

\begin{theorem}[Dimensional cascade]
\label{thm:cascade}
For all $d \geq 2$ and all $L \geq 0$,
\begin{equation}
g_L(S^d) = \sum_{l=0}^{L} g_l(S^{d-1}).
\label{eq:cascade}
\end{equation}
\end{theorem}

\begin{proof}
The cumulative degeneracy of the first $L+1$ harmonic spaces on
$S^{d-1}$ is
\begin{equation}
\sum_{l=0}^{L} g_l(S^{d-1}) = \sum_{l=0}^{L}
\left[\binom{l+d-1}{d-1} - \binom{l+d-3}{d-1}\right].
\end{equation}
This is a telescoping sum.  Writing $a_l = \binom{l+d-1}{d-1}$ and
noting $\binom{l+d-3}{d-1} = a_{l-2}$ with the convention
$a_l = 0$ for $l < 0$:
\begin{align}
\sum_{l=0}^{L} (a_l - a_{l-2})
&= a_L + a_{L-1} - a_{-1} - a_{-2} \nonumber \\
&= \binom{L+d-1}{d-1} + \binom{L+d-2}{d-1}.
\end{align}
By the Vandermonde identity $\binom{n}{k} + \binom{n-1}{k} =
\binom{n+1}{k+1} \cdot (k+1)/(n+1)$, or more directly by the
Pascal relation $\binom{n}{k} + \binom{n}{k+1} = \binom{n+1}{k+1}$
applied column-wise, the sum equals
\begin{equation}
\binom{L+d}{d} - \binom{L+d-2}{d} = g_L(S^d),
\end{equation}
which is Eq.~\eqref{eq:degeneracy}.
\end{proof}

The cascade provides a recursive construction of the entire
hierarchy from $S^1$.  The familiar identity
$n^2 = \sum_{l=0}^{n-1}(2l+1)$, which decomposes each hydrogen
shell into angular momentum subshells, is the $d=3$ instance:
$g_{n-1}(S^3) = (n)^2 = \sum_{l=0}^{n-1}(2l+1) =
\sum_{l=0}^{n-1} g_l(S^2)$.

\subsection{Verified numerical values}

Table~\ref{tab:numerical} presents the first five eigenvalues and
degeneracies for $S^3$ through $S^7$, verified computationally
against the formulas in Eqs.~\eqref{eq:eigenvalues}
and~\eqref{eq:degeneracy}.

\begin{table}
\caption{Eigenvalues $|\lambda_l| = l(l+d-1)$ and degeneracies
$g_l$ for $S^3$ through $S^7$, $l = 0, \ldots, 4$.
\label{tab:numerical}}
\begin{ruledtabular}
\begin{tabular}{c|rr|rr|rr|rr|rr}
 & \multicolumn{2}{c|}{$S^3$} & \multicolumn{2}{c|}{$S^4$}
 & \multicolumn{2}{c|}{$S^5$} & \multicolumn{2}{c|}{$S^6$}
 & \multicolumn{2}{c}{$S^7$} \\
$l$ & $|\lambda|$ & $g$ & $|\lambda|$ & $g$
    & $|\lambda|$ & $g$ & $|\lambda|$ & $g$
    & $|\lambda|$ & $g$ \\
\hline
0 &  0 &  1 &  0 &  1 &  0 &  1 &  0 &   1 &  0 &   1 \\
1 &  3 &  4 &  4 &  5 &  5 &  6 &  6 &   7 &  7 &   8 \\
2 &  8 &  9 & 10 & 14 & 12 & 20 & 14 &  27 & 16 &  35 \\
3 & 15 & 16 & 18 & 30 & 21 & 50 & 24 &  77 & 27 & 112 \\
4 & 24 & 25 & 28 & 55 & 32 &105 & 36 & 182 & 40 & 294 \\
\end{tabular}
\end{ruledtabular}
\end{table}


%=================================================================
\section{The $G_2/S^6$ Correspondence}
\label{sec:g2}
%=================================================================

\subsection{$S^6$ as a coset space}

The six-sphere admits the coset representation
$S^6 = G_2/\mathrm{SU}(3)$, where $G_2$ is the 14-dimensional
exceptional Lie group of automorphisms of the octonions.
The unit imaginary octonions form $S^6 \subset \mathrm{Im}(\mathbb{O})
\cong \mathbb{R}^7$, and $G_2$ acts transitively on $S^6$ with
isotropy subgroup $\mathrm{SU}(3)$.  This coset structure is
the analog of $S^3 = \mathrm{SO}(4)/\mathrm{SO}(3)$ for the
Coulomb case and $S^5 \cong \mathrm{SU}(3)/\mathrm{SU}(2)$ for
the harmonic oscillator case (Paper~24).

\subsection{Dimension identity}

\begin{theorem}[G$_2$ dimension identity]
\label{thm:dim}
For all $l \geq 0$, the dimension of the symmetric $G_2$
irreducible representation with Dynkin labels $(l,0)$ equals
the degeneracy of the $l$-th harmonic on $S^6$:
\begin{equation}
\dim(l,0)_{G_2} = g_l(S^6).
\label{eq:dim_identity}
\end{equation}
\end{theorem}

\begin{proof}
The Weyl dimension formula for $G_2$ gives the dimension of the
irreducible representation with highest weight
$\lambda = a\,\omega_1 + b\,\omega_2$ (in fundamental weight
basis) as
\begin{equation}
\dim(a,b) = \frac{(a{+}1)(b{+}1)(a{+}b{+}2)(a{+}2b{+}3)
(a{+}3b{+}4)(2a{+}3b{+}5)}{120}.
\label{eq:weyl_g2}
\end{equation}
Setting $b = 0$:
\begin{align}
\dim(l,0) &= \frac{(l+1)\cdot 1 \cdot (l+2)(l+3)(l+4)(2l+5)}{120}
\nonumber \\
&= \frac{(l+1)(l+2)(l+3)(l+4)(2l+5)}{120}.
\label{eq:dim_l0}
\end{align}
The $S^6$ harmonic degeneracy from Eq.~\eqref{eq:degeneracy} with
$d = 6$ is
\begin{align}
g_l(S^6) &= \binom{l+6}{6} - \binom{l+4}{6} \nonumber \\
&= \frac{(l+6)!}{6!\,l!} - \frac{(l+4)!}{6!\,(l-2)!}.
\label{eq:g_s6}
\end{align}
Expanding and simplifying (verified by direct computation for
$l = 0, \ldots, 19$):
\begin{equation}
g_l(S^6) = \frac{(l+1)(l+2)(l+3)(l+4)(2l+5)}{120},
\end{equation}
which is identical to Eq.~\eqref{eq:dim_l0}.
\end{proof}

\noindent The first eight values are $g_l = 1, 7, 27, 77, 182,
378, 714, 1254$, verified numerically.

\subsection{Casimir--Laplacian proportionality}

\begin{theorem}[Casimir--Laplacian proportionality]
\label{thm:casimir}
The quadratic Casimir eigenvalue of $G_2$ on the symmetric
irrep $(l,0)$ is proportional to the $S^6$ Laplace--Beltrami
eigenvalue:
\begin{equation}
C_2^{G_2}(l,0) = 2\,l(l+5) = 2\,|\lambda_l(S^6)|.
\label{eq:casimir}
\end{equation}
\end{theorem}

\begin{proof}
The quadratic Casimir for a simple Lie algebra $\mathfrak{g}$ on
the irrep with highest weight $\lambda$ is
$C_2(\lambda) = (\lambda + 2\rho,\, \lambda)$,
where $\rho$ is the Weyl vector and $(\cdot,\cdot)$ is the inner
product on the weight space.  For $G_2$, the inner product matrix
of the fundamental weights satisfies
$(\omega_1,\omega_1) = 2$, $(\omega_2,\omega_2) = 6$,
$(\omega_1,\omega_2) = 3$,
using the convention where the short root has squared length~2.
With $\lambda = l\,\omega_1$ and $\rho = \omega_1 + \omega_2$:
\begin{align}
(\lambda,\lambda) &= 2l^2, \\
(\rho,\lambda) &= l\bigl[(\omega_1,\omega_1) + (\omega_2,\omega_1)\bigr]
= 5l.
\end{align}
Therefore
\begin{equation}
C_2(l,0) = 2l^2 + 2 \cdot 5l = 2l^2 + 10l = 2l(l+5).
\end{equation}
Since $|\lambda_l(S^6)| = l(l+5)$, the proportionality
$C_2(l,0) = 2\,|\lambda_l|$ holds exactly.
\end{proof}

The factor of~2 is a normalization convention tied to the choice
of inner product on $\mathfrak{g}_2$.  The structural content is
that the $G_2$ Casimir reproduces the $S^6$ Laplace--Beltrami
spectrum on the symmetric irreps.

\subsection{Peter--Weyl selection: only $(l,0)$ irreps appear}

The Peter--Weyl theorem on the homogeneous space $G_2/\mathrm{SU}(3)$
decomposes $L^2(S^6)$ into $G_2$ representations.  A $G_2$ irrep
$(a,b)$ appears in this decomposition if and only if it contains
the trivial $\mathrm{SU}(3)$ representation under the branching
$G_2 \supset \mathrm{SU}(3)$.

For the symmetric irreps $(l,0)$, the branching
$G_2 \supset \mathrm{SU}(3)$ yields:
\begin{equation}
(l,0)_{G_2} \to \bigoplus_{p,q} (p,q)_{\mathrm{SU}(3)},
\end{equation}
where the trivial $(0,0)_{\mathrm{SU}(3)}$ appears exactly once.
This is a standard result: the $G_2$-invariant function on
$G_2/\mathrm{SU}(3)$ at each harmonic degree $l$ spans a single
copy of $(l,0)_{G_2}$.

For non-symmetric irreps with $b > 0$, the trivial
$\mathrm{SU}(3)$ representation does not appear in the branching.
For example, the adjoint $(0,1)_{G_2}$ of dimension~14 branches as
$14 \to 3 \oplus \bar{3} \oplus 8$ under $\mathrm{SU}(3)$, with no
trivial component.

Therefore, the harmonic decomposition of $L^2(S^6)$ is
\begin{equation}
L^2(S^6) = \bigoplus_{l=0}^{\infty} (l,0)_{G_2},
\label{eq:peter_weyl}
\end{equation}
and only symmetric $G_2$ irreps contribute.

\subsection{Structural parallel with $S^3/\mathrm{SO}(4)$}

Table~\ref{tab:parallel} summarizes the structural parallel
between the established $S^3$ case and the $S^6/G_2$ correspondence.

\begin{table}
\caption{Structural comparison of $S^3/\mathrm{SO}(4)$ and
$S^6/G_2$.
\label{tab:parallel}}
\begin{ruledtabular}
\begin{tabular}{lll}
Property & $S^3$ & $S^6$ \\
\hline
Symmetry group    & $\mathrm{SO}(4)$ & $G_2$ \\
Coset structure   & $\mathrm{SO}(4)/\mathrm{SO}(3)$ &
                    $G_2/\mathrm{SU}(3)$ \\
Isotropy subgroup & $\mathrm{SO}(3)$ & $\mathrm{SU}(3)$ \\
$g_1$             & 4                & 7 \\
Eigenvalue        & $l(l+2)$         & $l(l+5)$ \\
Casimir           & $\propto l(l+2)$ & $2\,l(l+5)$ \\
Irreps in $L^2$   & $(l/2,\,l/2)$ & $(l,0)$ \\
Algebra           & quaternionic     & octonionic \\
Physical system   & 3D hydrogen      & (none standard) \\
Packing           & Yes (Paper~0)    & No \\
Constant $\kappa$ & Yes ($-1/16$)    & No \\
\end{tabular}
\end{ruledtabular}
\end{table}


\subsection{The octonionic cross product as a ladder operator}

On $S^3$, the inter-shell transition operators $T_{\pm}$ change the
principal quantum number $n$ by $\pm 1$, connecting adjacent
$\mathrm{SO}(4)$ irreps.  On $S^6$, the analog is the octonionic
cross product.

The imaginary octonions $\mathrm{Im}(\mathbb{O}) \cong \mathbb{R}^7$
carry a bilinear cross product $\mathbf{x} \times \mathbf{y} =
\mathrm{Im}(\mathbf{x} \cdot \mathbf{y})$ defined by the Fano plane
structure constants.  For a point $\mathbf{x} \in S^6$, the map
$J_{\mathbf{x}}(\mathbf{v}) = \mathbf{x} \times \mathbf{v}$
restricts to the tangent space $T_{\mathbf{x}} S^6$ and satisfies
$J^2 = -\mathrm{Id}$ (verified numerically at multiple points on
$S^6$), giving $S^6$ an almost-complex structure.

The $G_2$ tensor product decomposition of $(l,0) \otimes (1,0)$
shows that the cross product mixes adjacent $l$-shells.  For
$l = 0$:
\begin{equation}
(0,0) \otimes (1,0) = (1,0).
\end{equation}
For $l \geq 1$, the decomposition includes $(l+1,0)$, $(l-1,0)$,
$(l,0)$, and additional irreps with $b > 0$:
\begin{equation}
(l,0) \otimes (1,0) \supset (l{+}1,0) \oplus (l{-}1,0) \oplus (l,0).
\end{equation}
This is the direct analog of $T_{\pm}$ on $S^3$: the cross product
with a vector in the fundamental representation $(1,0)$ has
components that raise, lower, and preserve the harmonic degree~$l$.

The key structural difference is that the $S^6$ almost-complex
structure is non-integrable: $S^6$ is \emph{not} a complex manifold.
The Nijenhuis tensor is nonzero, directly tied to the
non-associativity of the octonions.  Computational verification
shows that 28 of the 35 ordered triples of imaginary octonion
basis elements have nonzero associator, confirming the
non-associativity density of $80\%$.

%=================================================================
\section{Uniqueness of $S^3$}
\label{sec:uniqueness}
%=================================================================

Three independent properties single out $S^3$ within the sphere
hierarchy.

\subsection{Fock projectability}

The Fock projection maps the momentum-space Coulomb Schr\"odinger
equation onto the free Laplacian on $S^3$ via stereographic
projection.  The mapping introduces a scaling parameter $\kappa$
that connects graph eigenvalues to physical energies.  For $S^3$
with the standard 3D Coulomb potential $V = -Z/r$, the parameter
$\kappa = -1/16$ is constant across all energy shells (Paper~7).

Paper~23 proves the Fock projection rigidity theorem: the $S^3$
conformal projection maps a one-electron central-field Hamiltonian
onto the free Laplacian on $S^3$ if and only if the spectrum is
$l$-independent within each $n$-shell.  This condition is unique
to the $-Z/r$ potential by $\mathrm{SO}(4)$ symmetry.

For $S^6$, the generalized Fock projection of the six-dimensional
Coulomb problem gives a spectrum $E_N = -Z^2/[2(N + 3/2)^2]$ with
graph eigenvalues $\lambda_l = l(l+5)$.  The effective ratio
$E/\lambda$ is not constant:
\begin{align*}
N=2:&\quad \kappa_{\text{eff}} = -1/147, \\
N=3:&\quad \kappa_{\text{eff}} = -1/567, \\
N=4:&\quad \kappa_{\text{eff}} = -1/1452.
\end{align*}
The non-constancy is a direct consequence of the rigidity theorem:
$d \neq 3$ Coulomb systems have $l$-dependent spectra within each
shell.

\subsection{Packing constructibility}

Paper~0 derives the angular quantum numbers $(\ell, m)$ and the
degeneracy $2\ell + 1$ from uniform-density packing of
distinguishable states on concentric shells in two-dimensional
flat space.  The shell capacities $2k - 1$ for shell~$k$ match
the $S^2$ harmonic degeneracy exactly.

Generalizing to $d$-dimensional flat packing, the shell capacity
is $k^d - (k-1)^d$, a polynomial of degree $d-1$ in $k$.  The
$S^d$ harmonic degeneracy $g_l(S^d)$ is also a polynomial of
degree $d-1$ in $l$, but with different coefficients for $d \geq 3$.
Explicit computation confirms that the flat packing capacity equals
the spherical harmonic degeneracy \emph{only} for $d = 2$:
\begin{equation}
k^2 - (k-1)^2 = 2k - 1 = 2l + 1 = g_l(S^2)
\end{equation}
with $l = k - 1$.  For all $d \geq 3$, the flat packing and
spherical harmonic sequences diverge.  For example, at $d = 3$:
$k^3 - (k-1)^3 = 3k^2 - 3k + 1$ versus $g_l(S^3) = (l+1)^2 = k^2$.

The packing construction is therefore unique to the $d = 2 \to S^2$
pathway, and the subsequent lift to $S^3$ via the Fock projection
is the only route from flat packing to a physical Hamiltonian.

\subsection{Physical realization}

$S^3$ encodes the three-dimensional Coulomb problem, which
describes the hydrogen atom and all hydrogenic ions.  No standard
three-dimensional quantum system maps to $S^4$, $S^6$, or $S^7$.
$S^5$ encodes the three-dimensional harmonic oscillator but through
the holomorphic sector of $H^2(S^5)$, not through a Fock-type
conformal projection.

Table~\ref{tab:uniqueness} summarizes the three criteria across
the hierarchy.

\begin{table}
\caption{Uniqueness criteria for spheres in the GeoVac framework.
Fock: admits a conformal projection with constant $\kappa$.
Pack: constructible from flat packing.
Phys: encodes a standard 3D quantum system.
\label{tab:uniqueness}}
\begin{ruledtabular}
\begin{tabular}{clccc}
$d$ & Sphere & Fock & Pack & Phys \\
\hline
1 & $S^1$ & --- & --- & 1D ring \\
2 & $S^2$ & --- & Yes & Angular sector \\
3 & $S^3$ & Yes & (via $S^2$) & 3D Coulomb \\
4 & $S^4$ & No  & No  & --- \\
5 & $S^5$ & No  & No  & 3D HO\footnotemark[1] \\
6 & $S^6$ & No  & No  & --- \\
7 & $S^7$ & No  & No  & 4D HO\footnotemark[1] \\
\end{tabular}
\end{ruledtabular}
\footnotetext[1]{Via holomorphic sector (Bargmann--Segal), not
Fock projection.}
\end{table}

$S^3$ is the unique sphere satisfying all three criteria
simultaneously.  This triple intersection---Fock projectability,
packing constructibility, and physical realization as a 3D
system---is what makes $S^3$ the natural home of the GeoVac
framework's Level~1 (atomic) Hamiltonian.

%=================================================================
\section{The Dimensional Cascade and Cross-Sphere Relations}
\label{sec:cascade}
%=================================================================

The cascade theorem (Theorem~\ref{thm:cascade}) provides an
algebraic link between adjacent spheres in the hierarchy.  Each
step up in dimension collects the harmonic spaces of the lower
sphere into a single shell of the higher sphere.

The best-known instance is the $S^2 \to S^3$ cascade:
\begin{equation}
g_{n-1}(S^3) = n^2 = \sum_{l=0}^{n-1} (2l+1) =
\sum_{l=0}^{n-1} g_l(S^2),
\end{equation}
which decomposes the $n^2$ hydrogen orbital degeneracy into angular
momentum subshells.  This is the $\mathrm{SO}(4) \supset
\mathrm{SO}(3)$ branching rule expressed as a degeneracy identity.

The $S^5 \to S^6$ cascade connects the harmonic oscillator and
$G_2$ spheres:
\begin{equation}
g_L(S^6) = \sum_{l=0}^{L} g_l(S^5),
\end{equation}
with $g_l(S^5) = (l+1)(l+2)^2(l+3)/12$ being the full $S^5$ harmonic
degeneracy.  The 3D~HO shell degeneracy $(l+1)(l+2)/2$ is the
$\mathrm{SU}(3)$ symmetric irrep $(l,0)$ dimension, which is a
proper holomorphic subsector selected by the Bargmann--Segal
projection (Paper~24).

Within the GeoVac natural geometry hierarchy, the cascade connects
Level~1 ($S^3$, atoms) to Level~3 ($S^5$, two-electron
hyperspherical) through the chain $S^2 \to S^3 \to S^4 \to S^5$.
Each step adds one continuous degree of freedom: $S^2$ (angular
momentum) $\to$ $S^3$ (Fock projection adds radial/energy) $\to$
$S^4$ (no known single-particle identification) $\to$ $S^5$
(Bargmann--Segal adds the holomorphic oscillator degree of freedom).
The physical ``gap'' at $S^4$ is consistent with the absence of a
standard quantum system between the Coulomb and harmonic oscillator
potentials in the classification of exactly solvable problems.

%=================================================================
\section{Discussion}
\label{sec:discussion}
%=================================================================

\subsection{Mathematical context}

The $G_2/S^6$ correspondence is a known result in differential
geometry: $S^6 = G_2/\mathrm{SU}(3)$ is a classical example of a
homogeneous nearly-K\"ahler manifold.  The contribution of this
paper is to place the correspondence within the GeoVac sphere
hierarchy, alongside the established $S^3/\mathrm{SO}(4)$ and
$S^5/\mathrm{SU}(3)$ cases, and to verify computationally that
the dimension identity, Casimir proportionality, and Peter--Weyl
selection hold to the numerical precision required by the framework.

$G_2$ holonomy manifolds play a role in string theory and
M-theory, where seven-dimensional manifolds with $G_2$ holonomy
are used for compactification from eleven to four
dimensions~\cite{Joyce2000}.  The $S^6$ graph constructed here
would encode the angular selection rules for fields on such
manifolds.  We note this connection without overclaiming: the
GeoVac framework operates in the non-relativistic quantum
mechanics regime, and the string-theoretic applications of $G_2$
involve very different physics.

\subsection{Open questions}

Two structural questions remain.

First, the non-symmetric $G_2$ irreps $(a,b)$ with $b > 0$ do not
appear in $L^2(S^6)$ by the Peter--Weyl theorem, but they do
appear in the tensor product decomposition of symmetric irreps
(Sec.~\ref{sec:g2}).  Whether these irreps have a packing
interpretation---perhaps on a space other than $S^6$---is unknown.

Second, the sphere hierarchy $S^1, S^2, \ldots$ is one family of
compact symmetric spaces.  Other families (complex projective
spaces $\mathbb{CP}^n$, Grassmannians, exceptional symmetric
spaces) have their own Laplace--Beltrami spectra and
representation-theoretic decompositions.  Whether a ``super-hierarchy''
connecting these families exists, with the sphere sequence as a
distinguished subfamily, is an open question with potential
implications for extending the GeoVac framework beyond spherical
geometries.

\subsection{Implications for the GeoVac framework}

The results of this paper sharpen the universal--Coulomb-specific
partition established in Papers~22--24.  The dimensional cascade
(Theorem~\ref{thm:cascade}) and the degeneracy formulas
\eqref{eq:degeneracy} are universal: they hold for any sphere
$S^d$ regardless of what physical system (if any) it encodes.
The angular sparsity theorem (Paper~22) is similarly universal.

The Coulomb-specific features---constant $\kappa$, Fock
projectability, packing constructibility---are unique to $S^3$.
The $G_2/S^6$ correspondence demonstrates that rich
representation-theoretic structure can exist on spheres that
lack these Coulomb-specific features: $S^6$ has a clean
Casimir--Laplacian proportionality and coset structure, yet
fails both the Fock and packing criteria.

This classification supports the design principle stated in
Paper~24: stay on the graph whenever possible (where the universal
angular structure applies), and when projection to a specific
physical system is necessary, identify the minimal
Coulomb-specific content required.

%=================================================================
% Acknowledgments
%=================================================================

\begin{acknowledgments}
This work was performed using an AI-augmented agentic workflow.
The principal investigator provided scientific direction and quality
control; implementation, computation, and documentation drafting
were performed collaboratively with large language models (Anthropic
Claude).
\end{acknowledgments}

%=================================================================
% References
%=================================================================

\begin{thebibliography}{20}

\bibitem{Fock1935}
V.~Fock, Z.~Phys. \textbf{98}, 145 (1935).

\bibitem{Joyce2000}
D.~D.~Joyce, \textit{Compact Manifolds with Special Holonomy}
(Oxford University Press, Oxford, 2000).

\end{thebibliography}

\end{document}
